\chapter{كيف يدخل العالم}
% full version: chapters/11_sensors.tex

كل عضو حس محوّل، كالميكروفون أو الخلية الضوئية. الضوء أو الصوت أو الجزيء يؤثر في بروتين خاص في الخلية الحساسة، فيفتح البروتين ”صنبورًا“، فيسري تيار، وفي آخر السلسلة تخرج نقرات تذهب إلى الشبكة. كل المستشعرات تقريبًا تقذف الغلوتامات، فتتصل مداخل الآلة مباشرة بشبكتها الهاتفية.

\section*{عند حدود الفيزياء}

تعمل مستشعرات الدماغ عند أقصى حدود الممكن. مستشعر الضوء في الظلام يستجيب لجسيم ضوء واحد. ومستشعر الصوت يلحظ إزاحة أقل من نانومتر، أي أقل من حجم بضع ذرات. وفي العتمة لا تحدّ الدقةَ العينُ، بل الضوء نفسه: الجسيمات تصل عشوائيًا، ولتمييز درجة اللون يجب جمع المزيد منها. لذلك نرى في الغسق أبطأ: تجمع العين الضوء مدة أطول، كآلة تصوير بزمن تعريض طويل.

\section*{التعريض التلقائي}

تتغير الإضاءة من ليلة مرصعة بالنجوم إلى ثلج تحت الشمس مليارات المرات، ويتغير علوّ الصوت أكثر من ذلك. أما تردد النقرات فيتغير من الصفر إلى بضع مئات في الثانية. لا يمكن حشر الأول في الثاني مباشرة. والحل هو حل آلة التصوير نفسه: التعريض التلقائي. يضبط المستشعر حساسيته طوال الوقت على المستوى المتوسط من حوله، ولا يبلغ عن السطوع، بل عن السطوع \emph{مقارنة بالخلفية}. الثابت يكف مع الوقت عن إثارة الاستجابة، وكل تغيير يثيرها. مداخل الآلة تتحدث عن التغيرات منذ البداية.

\pencilfig{125}{%
% light rises in two tenfold steps, then drops a hundredfold; sensor rate = baseline + gain * (log light - adapted level),
% the adapted level follows log light with a 0.7 s time constant; rate clipped at zero
\draw[penlite=pcAmber] (0,1.75) -- (1.4,1.75) -- (1.4,2.1) -- (4.2,2.1) -- (4.2,2.45) -- (6.3,2.45) -- (6.3,1.75) -- (8.4,1.75);
\node[lbl, text=pcAmber!70!black, anchor=east] at (-0.05,2.0) {الضوء};
\node[lbl, anchor=south] at (2.8,2.12) {أسطع 10 مرات}; \node[lbl, anchor=south] at (5.25,2.47) {ثم 10 أخرى};
\node[lbl, anchor=south] at (7.35,1.77) {عتمة من جديد};
\draw[pen] (0,0) -- (8.4,0);
\draw[pcGray, densely dashed, line width=0.5pt] (0,0.32) -- (8.4,0.32);
\draw[penlite=pcBlue] plot coordinates {(0.00,0.32) (1.26,0.32) (1.40,1.02) (1.54,0.85) (1.68,0.72) (1.82,0.62) (1.96,0.54) (2.10,0.49) (2.24,0.45) (2.38,0.41) (2.52,0.39) (2.66,0.37) (2.80,0.36) (3.08,0.34) (3.50,0.33) (4.06,0.32) (4.20,1.03) (4.34,0.85) (4.48,0.72) (4.62,0.62) (4.76,0.54) (4.90,0.49) (5.04,0.45) (5.18,0.41) (5.32,0.39) (5.46,0.37) (5.60,0.36) (5.88,0.34) (6.16,0.33) (6.30,0.00) (7.00,0.00) (7.14,0.07) (7.28,0.13) (7.42,0.18) (7.56,0.22) (7.70,0.24) (7.84,0.26) (7.98,0.28) (8.12,0.29) (8.40,0.30)};
\node[lbl, text=pcBlue, anchor=east] at (-0.05,0.32) {الاستجابة};
\node[lbl, anchor=north] at (6.65,-0.02) {صامت};
}{يزداد الضوء سطوعًا على درجات، كل منها عشرة أضعاف. يستجيب المستشعر لكل تغيير بقفزة ويعود سريعًا إلى مستواه السابق. وحين تعود العتمة، يصمت مدة.}

كرر المهندسون هذه الحيلة في كاميرات الأحداث. كاميرا كهذه لا تلتقط الإطارات وفق ساعة: كل بكسل فيها يبلغ بنفسه ”صار أسطع“ أو ”صار أعتم“ حين يتغير شيء. إنه زوج السلكين ”نعم“ و”لا“ نفسه من الفصل الثاني، لكن في السيليكون.

\section*{العين تضغط الصورة}

في العين نحو مئة مليون مستشعر ضوئي، أما الأسلاك التي تنقل الصورة إلى الدماغ فنحو مليون. قبل أن تغادر الصورة العين تُضغط نحو مئة مرة. المساحات المتجانسة لا تكلف شيئًا تقريبًا، وتُنقل الحواف والتغيرات. ووفق تقدير أُجري على عيون الحيوانات ثم حُوّل إلى الإنسان، تنقل العين إلى الدماغ نحو عشرة ملايين بت في الثانية، ككابل شبكة قديم.

\section*{الأذن بيانو}

تحلل الأذن الصوت إلى تردداته، كما كان المهندس سيضع مجموعة من المرشحات. غشاء مرن في الأذن يستجيب في مواضعه المختلفة لترددات مختلفة، والمستشعرات في كل موضع تصغي إلى ”مفتاحها“. رقم المستشعر الذي عمل هو طبقة الصوت.

\pencilfig{11}{\pencilart{11}}{الأذن مبنية كبيانو معكوس: كل مفتاح يصغي إلى نغمته.}

بل إن الأذن ليست سلبية. بعض خلاياها يهز الغشاء بنفسه بإيقاع الصوت ويضخم الأصوات الخافتة آلاف المرات، ولا يكاد يضخم العالية. إنه مضخم عند حافة الإثارة الذاتية بالضبط: هناك يكون أشد حساسية. وعند الترددات المنخفضة تنقر العصبونات أيضًا بإيقاع الموجة، فتحفظ التوقيت الدقيق للصوت، ومنه يجد الدماغ اتجاه المصدر. وكلما ارتفع التردد، عجزت العصبونات أكثر عن مجاراته، ولا يبقى إلا الترميز بالموضع.

\section*{سائر الحواس}

في كل حاسة تقريبًا نتعرف على حيلة من الفصول السابقة. الحرارة ينقلها سلكان، واحد للدفء وآخر للبرد. واللمس يستخدم مستشعرات تستجيب للمسة ثم تصمت فورًا، ومستشعرات تحافظ على استجابتها ما دام الضغط قائمًا. أما الشم فمبني كالرمز الشريطي: لدى الإنسان نحو أربعمئة نوع من المستشعرات الشمية، والرائحة هي مجموعة الأنواع التي عملت. وعدد هذه المجموعات فلكي.

\fullversion{11}
